% ============================================================
% OP_GUT2_repclass_recon_v2.tex  (PROPOSAL-ONLY; v2 per GPT review)
% Changes vs v1: "P7-restricted" rename; structural-restriction and
% no-exotics caveats; minimality formally defined as a convention;
% LH-Weyl conventions fixed; full anomaly system incl. SU(3)^3
% written out; SM-pattern sector solved by reuse of recorded
% AN3/N_c algebra; nu^c three-branch treatment consistent with
% AN3; expectation language removed from solution-space task;
% structural lemmas proven (vector-like transparency; no chiral
% pure-singlet sets with <=4 fields; <=2 coloured weak-singlets
% forced vector-like; colourless doublet sector closed at minimal
% multiplicity); case tree with computed vs explicitly-open
% sectors; status split candidate-Level-A-algebra / Level-B-ECT /
% OP-GUT2-Open; self-contained (no citations). NO preprint edits.
% ============================================================
\section*{OP-GUT2 conditional solution-space classification, v2
(proposal-only)}

\textit{Status target: candidate Level~A algebra inside the
restricted class (for the sectors actually computed); Level~B
conditional as ECT application, because the admissible class is
structurally imposed; OP-GUT2 itself remains Open.}

\subsection*{1. Conventions}
All fermion fields are written as left-handed Weyl fields;
$u^c,d^c,e^c$ (and, in its branch, $\nu^c$) denote left-handed
conjugates of the right-handed Standard-Model fields.
Anomaly coefficients are normalised so that
$A(\mathbf 3)=+1$, $A(\bar{\mathbf 3})=-1$ for the cubic colour
anomaly; $\mathbf 2$ of $SU(2)$ is pseudo-real, so
$(\mathbf 1,\mathbf 2)_Y\oplus(\mathbf 1,\mathbf 2)_{-Y}$ is a
vector-like (Dirac-pairable) combination.

\subsection*{2. Definition: P7-restricted admissible matter class}
A \emph{P7-restricted admissible matter class} element is a
left-handed Weyl field valued in a tensor product of the already
assumed structures only: the colour module $E$ or $E^*$
($\mathbf 3$ or $\bar{\mathbf 3}$), the weak doublet/singlet
structure ($\mathbf 2$ or $\mathbf 1$), and a compact-$U(1)$
charge-lattice label $Y$.
This is a structural restriction defining a test class, not a
derivation of matter representations.
The exclusion of higher colour tensors
($\mathbf 6,\mathbf 8,\dots$) and higher weak isospin is not
derived from anomaly cancellation and not derived from ECT
dynamics; it is a minimality/available-module assumption used to
define the test class.

\subsection*{3. Minimality (classification convention)}
A \emph{chiral type} is a triple $(R_c,R_w,Y)$ with
$R_c\in\{\mathbf 1,\mathbf 3,\bar{\mathbf 3}\}$,
$R_w\in\{\mathbf 1,\mathbf 2\}$, $Y$ on the lattice.
\emph{Minimality} means: at most one copy of each chiral type;
no vector-like pairs counted as part of the minimal chiral core
(a vector-like pair is $(R_c,R_w,Y)\oplus(\bar R_c,R_w,-Y)$, with
$R_w=\mathbf 2$ self-conjugate); and no sterile singlets unless an
explicit $\nu^c$ branch is being analysed.
This is a classification convention, not a dynamical selection
rule.

\subsection*{4. Anomaly system on the class}
For a set $\{(R_{c,i},R_{w,i},Y_i)\}$ with colour/weak dimensions
$d_{c,i},d_{w,i}$:
\begin{align*}
&[SU(3)]^3: && \textstyle\sum_i d_{w,i}\,A(R_{c,i})=0;\\
&[SU(3)]^2U(1): && \textstyle\sum_{i\,\rm coloured} d_{w,i}\,Y_i=0;\\
&[SU(2)]^2U(1): && \textstyle\sum_{i\,\rm doublets} d_{c,i}\,Y_i=0;\\
&\text{grav}^2\,U(1): && \textstyle\sum_i d_{c,i}d_{w,i}\,Y_i=0;\\
&[U(1)]^3: && \textstyle\sum_i d_{c,i}d_{w,i}\,Y_i^3=0;\\
&\text{Witten}: && \textstyle\sum_{i\,\rm doublets} d_{c,i}
  \equiv 0 \pmod 2 .
\end{align*}
The Witten count is colour-weighted: $(\mathbf 3,\mathbf 2)$
contributes three doublets.

\subsection*{5. SM-pattern sector (computed; reuses recorded
AN3/$N_c$ algebra)}
Pattern $\mathcal P_{\rm SM}=
\{(\mathbf 3,\mathbf 2)_{Y_Q},(\bar{\mathbf 3},\mathbf 1)_{Y_u},
(\bar{\mathbf 3},\mathbf 1)_{Y_d},(\mathbf 1,\mathbf 2)_{Y_L},
(\mathbf 1,\mathbf 1)_{Y_e}\}$.
$[SU(3)]^3$ explicitly: $Q$ contributes two colour triplets
($d_w=2$), $u^c,d^c$ two anti-triplets, so
$2A(\mathbf 3)+2A(\bar{\mathbf 3})=2-1-1=0$.
Witten: $3+1=4$, even.
The remaining conditions are
$2Y_Q+Y_u+Y_d=0$, $3Y_Q+Y_L=0$,
$6Y_Q+3Y_u+3Y_d+2Y_L+Y_e=0$,
$6Y_Q^3+3Y_u^3+3Y_d^3+2Y_L^3+Y_e^3=0$ ---
exactly the system whose solution structure is already recorded
(AN3 cubic, its two hypercharge branches, the degenerate $Y_Q=0$
branch, and the overall normalisation freedom; the $N_c=3$ case of
the recorded $N_c$-classification).
Nothing new is asserted here; the SM one-generation assignments
appear as a solution branch.
\emph{$\nu^c$ branches (consistent with AN3):}
(a)~no $\nu^c$ (baseline above);
(b)~hypercharge-neutral singlet, $Y_{\nu^c}=0$: all conditions
unchanged;
(c)~free-hypercharge $\nu^c$: opens the recorded $B-L$ flat
direction.
$\nu^c$ is not part of the baseline set.

\subsection*{6. Structural lemmas (own algebra; proofs sketched,
to be written in full)}
\paragraph{Lemma V (vector-like transparency).}
A vector-like pair contributes zero to all five local conditions
and an even number to the Witten count.
Hence anomaly conditions do not constrain vector-like additions;
they are excluded from the minimal core by convention, not by
anomalies.
\paragraph{Lemma N$_{\le4}$ (no small chiral pure-singlet sets).}
Within the class, there is no anomaly-free chiral set consisting
solely of $n\le4$ fields $(\mathbf 1,\mathbf 1)_{y_i}$ with all
$y_i\neq0$.
Proof sketch: the conditions reduce to $\sum y_i=0$ and
$\sum y_i^3=0$.
$n=1$: $y=0$.
$n=2$: $y_2=-y_1$, a vector-like pair.
$n=3$: $\sum y_i=0$ implies $\sum y_i^3=3y_1y_2y_3$, so some
$y_i=0$.
$n=4$: write $y_1+y_2=s=-(y_3+y_4)$; then
$y_1^3+y_2^3=s^3-3y_1y_2s$ and
$(-y_3)^3+(-y_4)^3=s^3-3y_3y_4s$; the cubic condition forces
$s\,(y_1y_2-y_3y_4)=0$; $s=0$ gives a vector-like pair, otherwise
$\{y_1,y_2\}=\{-y_3,-y_4\}$, i.e.\ two vector-like pairs.
\paragraph{Lemma S$_2$ (small coloured weak-singlet sector).}
With no doublets and at most two coloured multiplets, $[SU(3)]^3$
forces one $\mathbf 3$ and one $\bar{\mathbf 3}$, and
$[SU(3)]^2U(1)$ forces $Y_{\bar{\mathbf 3}}=-Y_{\mathbf 3}$: a
vector-like pair.
\paragraph{Lemma D$_0$ (colourless doublet sector at minimal
multiplicity).}
Witten parity requires an even doublet count; two distinct
colourless doublets need $Y_1+Y_2=0$ by $[SU(2)]^2U(1)$, which by
pseudo-reality is a vector-like pair; a doublet at $Y=0$ is
excluded at odd count by Witten parity and duplicates are excluded
by the type convention.
Hence the colourless sector at minimal multiplicity reduces to the
pure-singlet problem of Lemma~N$_{\le4}$.

\subsection*{7. Solution-space classification task (no
expectations stated)}
Classify all non-SM chiral combinations inside the restricted
class that pass the anomaly conditions, including possible
vector-like additions, singlet branches, and $B-L$/hypercharge
freedoms.
No uniqueness claim is made before the solution space is
explicitly computed.
Organisation: a finite case tree over multiplet patterns with at
most six core multiplets (the SM-pattern size plus the $\nu^c$
branch), partitioned by coloured content and doublet count.
\emph{Computed so far:} the SM-pattern sector (\S5);
the sectors closed by Lemmas N$_{\le4}$, S$_2$, D$_0$.
\emph{Explicitly open cases (recorded, not guessed):}
pure-singlet sets with $\ge5$ fields (chiral solutions of the two
Diophantine conditions may exist; they carry no colour/weak
structure either way);
coloured weak-singlet sectors with $\ge3$ coloured multiplets;
mixed patterns at the multiplicity boundary (e.g.\ doublet sectors
combining $(\mathbf 3,\mathbf 2)$ with additional coloured
singlets beyond the SM pair).
The insertion, if approved later, will present computed sectors as
computed and open cases as open.

\subsection*{8. No-overclaim lemma (draft for the eventual
insertion)}
Anomaly cancellation, Witten parity, and lattice quantisation do
not by themselves select the Standard-Model matter content.
Within the P7-restricted class and under the stated minimality
convention, the one-generation SM pattern arises as a computed
solution sector; identifying it as the realised matter content
requires additional input (the structural class restriction
itself, the minimality convention, and ultimately the open
dynamical derivation OP-GUT2).

\subsection*{9. Literature policy}
Self-contained algebra first; no citations in the first round.
Geng--Marshak and Minahan--Ramond--Warner may later be added as
context only, after page-level source verification, and the
classification must not depend on borrowed uniqueness theorems
whose assumptions may differ from the P7-restricted class.

\subsection*{10. Questions for GPT}
(Q1)~Is the six-multiplet core bound an acceptable scope for the
case tree, with larger patterns recorded as out-of-scope?
(Q2)~Lemmas V, N$_{\le4}$, S$_2$, D$_0$: include in the eventual
insertion as candidate Level~A algebra inside the restricted
class, or keep N$_{\le4}$/D$_0$ proposal-only?
(Q3)~For the open coloured $\ge3$-multiplet sector: attempt full
computation in v3, or insert with the honest computed/open split
as in \S7?
(Q4)~After your review: v3 (if more algebra required) or exact
insertion blocks?
